\subsection{Petri nets}\label{sec:petri-nets}

Petri nets are a standard formalism for concurrency and process mining. We treat both classical Petri nets and object-centric Petri nets in a single typed framework. The common structural core is a finite place-transition incidence structure. The difference lies in the interpretation of tokens and markings. In the classical case, tokens are indistinguishable and markings count token multiplicity. In the object-centric case, places carry types and markings record finite collections of object identities of the appropriate type. In both settings, transitions will be interpreted as generator cospans of a cospan-algebra signature, and the flow between them as typed wires of its free symmetric monoidal category.

\begin{definition}[Object-centric Petri net~\cite{vandettenDiscoveringCompactLive2024}]
  An object-centric Petri net (OCPN) is a tuple $N=(P,T,F,\mathcal{O},\mathrm{type},\ell,M_0)$ of places $P$, transitions $T$, a flow relation $F\subseteq(P\times T)\cup(T\times P)$, a finite set of object types $\mathcal{O}$, a typing $\mathrm{type}:P\to\mathcal{O}$, a labelling $\ell:T\to\mathcal{L}\cup\{\tau\}$ over an activity alphabet $\mathcal{L}$ with silent symbol $\tau$, and an initial marking $M_0\in\mathbb{N}^{P}$. Write $T_\tau=\ell^{-1}(\tau)$ for the silent transitions and $\mathrm{src}(t),\mathrm{tgt}(t)$ for the input and output places of a transition $t$. Several transitions may share a label, the firing modes of one activity (Remark~\ref{rem:mode-refinement}).
\end{definition}

\begin{definition}[Petri net]
  A Petri net is an object-centric Petri net where $\mathrm{type}(p)=*$ for all $p\in{}P$.
\end{definition}

Both classes share this incidence structure and differ only in how tokens are read.
The classical case has indistinguishable counts, the object-centric case typed object
identities. One construction serves both. Before reading object types off the
net we record the structural condition that keeps them unambiguous across silent
transitions.

\begin{definition}[Type-consistent OCPN]\label{def:type-consistent}
  Let $(P\cup T_\tau,\,F)$ be the silent-flow subgraph of an OCPN,
  obtained by deleting the observable transitions $T\setminus T_\tau$ together
  with their incident arcs. The net is type-consistent when, along every
  directed path of the silent-flow subgraph between two places, the
  concretely-typed places carry one common object type. Untyped routing places
  and the artificial source and sink are exempt and inherit the type of the
  path on which they lie. Equivalently, no directed chain of silent transitions
  joins two places of different concrete object type. Exclusive routing is
  unaffected. A single untyped place may still fan out to alternatives of
  different type, since these lie on distinct directed paths.
\end{definition}

We assume throughout that object-centric nets are type-consistent. The
condition holds automatically for nets produced by object-centric discovery,
where each object type flows in its own coloured sub-net, and is the Petri-net
form of the prohibition, due to \citet{lissObjectCentricCausalNets2025} and reflected in the causal-net
boundary convention of Section~\ref{sec:causal-nets}, against a transition
that converts one object type into another. Remark~\ref{rem:ocpn-silent-typing}
records how the construction relies on it.

This construction instantiates the general framework of Section~\ref{sec:general-framework}.
Its content is the way we read a net as an AND/OR graph $(V,E)$. We build the mediators the
framework needs instead of reusing $P\cup T$ verbatim. Each place becomes an XOR mediator, since a
token in it is left by one of its producing transitions and taken by one of its consumers, a
choice. Each transition's firing becomes an AND synchronisation, consuming its whole pre-set and
yielding its whole post-set at once, so its input places are joined by an AND mediator and its
output places split from one. An observable transition is in addition the activity it labels,
sitting between its pre-set and post-set AND mediators. The activity does no combining of its own,
exactly as in Definition~\ref{def:lm-graph}, with all concurrency in the AND mediators
and all choice in the XOR places. A silent transition contributes only its transparent AND mediator
and no generator. These choices determine the instance map.
\medskip
\begin{center}
  \small\begin{tabularx}{\linewidth}{@{}lX@{}}
    \hline
    General framework & OCPN \\
    \hline
    Activities $\mathcal{A}$ & observable transitions $T\setminus T_\tau$ \\
    AND mediators & each transition's pre-set join and post-set split; silent transitions $T_\tau$ (transparent, no generator) \\
    XOR mediators & places $P$ \\
    Edges $E$ & the flow $F$, threaded through these mediators \\
    \hline
  \end{tabularx}
\end{center}
\medskip
The walk of Definition~\ref{def:contexts-general} runs on this graph unchanged. From the pre-set
AND mediator of an observable transition it passes through each input place, an XOR mediator that
keeps the reached sets of its producers separate, and through any silent transition, and it stops
at the first observable transitions. The post-set is read in the same way. The backward and
forward families hold one reached set per joint choice of observable neighbour, each endpoint read
through the labelling $\ell$, so transitions sharing an activity label contribute the same labelled
typed wire and are identified at composition. Each transition--context pair yields one generator
cospan (Definition~\ref{def:generator-cospan-general}), and no context is discarded.

\begin{theorem}[Canonical cospan-algebra presentation of typed Petri nets]
  \label{thm:petri-to-cospan-signature}
  Let \(N=(P,T,F,\mathcal{O},\mathrm{type},\ell,M_0)\) be an object-centric Petri net. Then \(N\) canonically determines a cospan-algebra signature \(\Sigma\), and hence the free symmetric monoidal category \(\mathbf{F}(\Sigma)\).
\end{theorem}

\begin{proof}
  The schema of Section~\ref{sec:notation-by-notation} applies at the instance map above,
  finiteness of $P$ and $T$ giving the conditions of Definition~\ref{def:lm-graph} and finite
  reach families (Definition~\ref{def:contexts-general}). The delta is typing across the silent fabric. The
  object-centric case runs over $\mathcal{O}$ with each edge typed by $\mathrm{type}$, and the
  type of each $\tau$-mediated thread is well defined precisely because $N$ is type-consistent
  (Definition~\ref{def:type-consistent}, Remark~\ref{rem:ocpn-silent-typing}).
\end{proof}

\begin{remark}[Relation to prior categorical presentations]\label{rem:smc-positioning}
  The result that Petri nets present free symmetric monoidal categories (SMCs) is due to
  \citet{meseguerPetriNetsAre1990}. The open-net cospan
  framework appears in \citet{baezOpenPetriNets2020}.
  Theorem~\ref{thm:petri-to-cospan-signature} refines these by constructing
  explicit context-indexed generator cospans. Each transition--context pair $(t,c)$
  yields a decorated cospan whose boundary data is read directly from the local
  incidence structure, equipping the SMC presentation with the generator-level
  information required by the cospan-algebra framework of
  Section~\ref{sec:hypergraphs}.
  In these terms the refinement makes the Meseguer and Montanari SMC
  presentation into an explicit generator-level transformation, positioned
  alongside the connector-algebra tradition in the related work of
  Section~\ref{sec:related-work}.
\end{remark}

\begin{corollary}[Untyped Petri nets as a special case]\label{cor:petri-untyped}
  Setting $\mathcal{O}=\{\ast\}$ and $\mathrm{type}(p)=\ast$ for all $p\in P$ in
  Theorem~\ref{thm:petri-to-cospan-signature} recovers the classical Petri-net
  presentation (Corollary~\ref{cor:untyped-general}), in which transitions induce generators,
  every typed wire carries the trivial type, and composites are the morphisms of the free
  symmetric monoidal category $\mathbf{F}(\Sigma)$ (Definition~\ref{def:free-category}), joined
  port to port along matching typed wires.
\end{corollary}

\begin{remark}[Silent transitions and type consistency]\label{rem:ocpn-silent-typing}
  The object-centric Petri-net literature attaches no typing of its own to a
  silent transition. None is needed here, because a silent transition reads its
  object types from the colours of its adjacent places,
  exactly as an observable transition does. Since silent transitions are
  transparent mediators (instance map above), a $\tau$-mediated typed wire
  $(x,w,y)$ carries the set $w$ of types on the thread it traces through
  the silent fabric, and type consistency (Definition~\ref{def:type-consistent})
  is precisely the condition that makes $w$ a single type. Under it the
  harvested boundary stays in bijection with the typed dependency arcs
  $D\subseteq T\times\mathcal{O}\times T$ of the object-centric causal net
  (Section~\ref{sec:causal-nets}). Transparent crossing introduces no new boundary
  types, and in particular no spurious duplication of a typed wire across the whole of
  $\mathcal{O}$.

  Dropping type consistency readmits exactly the wild behaviour it is there to
  tame. A type-converting silent transition would make a single thread enter as
  one concrete type and leave as another. The walk records both types in the set $w$
  of the typed wire through it, and the analysis reports a wire whose path carries more than one
  type (Section~\ref{subsec:analysis}). The construction therefore exposes the conversion
  instead of hiding it behind an artificial typed boundary. The
  only escape is the deliberately lossy projection to $\mathcal{O}=\{\ast\}$
  (Corollary~\ref{cor:petri-untyped}), which recovers the connectivity of $D$ but
  discards all typing. This is the Petri-net counterpart of the routing-only
  gateway discipline of Remark~\ref{rem:bpmn-typed-gateways}.
\end{remark}

\begin{remark}[Silent routing leaves the signature fixed]\label{rem:ab-place-vs-tau}
The sequence $a\,;\,b$ presented as one place $a\to p\to b$, and the same sequence with $p$ split by a silent transition into $a\to p_1$, $\tau$, $p_2\to b$, yield the identical signature. In both presentations activity $a$ has successor $b$ and activity $b$ has predecessor $a$. The silent transition passes transparently and contributes no generator, so the split place introduces no structural difference. The same argument absorbs an implicit place, a place duplicating a causal constraint that the net already carries. It contributes an endpoint the port set already holds, and ports are identified by producing activity, object type, and consuming activity, so the duplicate leaves the generator unchanged. This is the invariance the two Petri nets of Figure~\ref{fig:routing-collapse} exhibit. Note the scope. A redundant place carrying another object type is not absorbed, because it contributes a genuinely new typed port.
\end{remark}

\begin{remark}[Firing modes]\label{rem:mode-refinement}
  A single transition fires on its whole pre-set and yields its whole post-set, so it cannot itself
  choose among exclusive typed inputs or outputs. Such a choice is encoded the standard Petri-net
  way, by several transitions sharing one activity label, one per admissible firing mode, which the
  labelling $\ell$ re-identifies so that a downstream activity sees one labelled port rather than one
  per mode. This is the Petri-net analogue of the object-centric causal net's binding sets
  (Section~\ref{sec:causal-nets}). The same activity acquires one generator per mode, a larger
  signature that induces the same language (Section~\ref{sec:conversion-framework}).
\end{remark}

\begin{figure*}[tp]
  \centering
  \resizebox{\textwidth}{!}{\begin{tikzpicture}[petriBase, scale=0.95, inline,
    silentT/.style={transition, fill=black!22}]
  \node[T] (g1) at (0,-1.0625) {$\bot$};
  \node[blueP] (po) at (1.7,0) {};
  \node[T] (a)  at (3.4,0.85) {$a$};
  \node[silentT] (ta) at (3.4,-0.85) {$\tau$};
  \node[blueP] (pc) at (5.1,0) {};
  \node[T] (b)  at (6.7,0) {$b$};
  \node[blueP] (pr) at (8.4,0) {};
  \node[redP] (pi) at (3.4,-2.125) {};
  \node[T] (s) at (10.2,-1.0625) {$s$};
  \node[blueP] (pso) at (12.0,0) {};
  \node[redP] (psi) at (12.0,-2.125) {};
  \node[T] (g2) at (13.8,-1.0625) {$\top$};

  \draw[blueflow] ($(g1.north east)!0.25!(g1.south east)$) -- (po);
  \draw[redflow]  ($(g1.north east)!0.75!(g1.south east)$) -- (pi);
  \draw[blueflow] (po) -- (a);
  \draw[blueflow] (po) -- (ta);
  \draw[blueflow] (a) -- (pc);
  \draw[blueflow] (ta) -- (pc);
  \draw[blueflow] (pc) -- (b);
  \draw[blueflow] (b) -- (pr);
  \draw[blueflow] (pr) -- ($(s.north west)!0.25!(s.south west)$);
  \draw[redflow]  (pi) -- ($(s.north west)!0.75!(s.south west)$);
  \draw[blueflow] ($(s.north east)!0.25!(s.south east)$) -- (pso);
  \draw[redflow]  ($(s.north east)!0.75!(s.south east)$) -- (psi);
  \draw[blueflow] (pso) -- ($(g2.north west)!0.25!(g2.south west)$);
  \draw[redflow]  (psi) -- ($(g2.north west)!0.75!(g2.south west)$);

  \node[blueP, minimum size=5mm, label={[font=\small]right:order}]  at (15.3,-0.55) {};
  \node[redP,  minimum size=5mm, label={[font=\small]right:item}] at (15.3,-1.55) {};
\end{tikzpicture}
\unskip}
  \caption{A small object-centric Petri net. An order (blue) is registered by $b$,
    optionally after a check $a$ that may be skipped through the silent transition
    $\tau$ (drawn filled). The registered order and an item (red) are then shipped
    together by $s$.}
  \label{fig:re-ocpn}
\end{figure*}

\begin{figure*}[tp]
  \centering
  \begin{tikzpicture}[inline,
    bspd/.style={spider, fill=bluecol},
    rspd/.style={spider, fill=redcol},
    mbox/.style={draw=black, rounded corners=4pt, fill=white, minimum size=0.75cm, font=\small},
    plabel/.style={font=\tiny}]

  \node[font=\footnotesize] at (2.7,1.4) {decorated cospan};
  \node[font=\footnotesize] at (9.0,1.4) {string diagram};

  \node[font=\small] at (-2.15,0) {$g_s$};
  \node[cospanblue, minimum height=0.9cm] (sL) at (0,0) {};
  \node[bspd] at (0,0.20) {};
  \node[rspd] at (0,-0.20) {};
  \node at (0.9,0) {$\rightarrow$};
  \node[cospangrey, minimum height=1.0cm, minimum width=2.6cm] (sA) at (2.7,0) {};
  \node[mbox, minimum height=0.8cm] (sf) at (2.7,0) {$s$};
  \node[bspd] (sao) at (1.65,0.20) {};
  \node[rspd] (sai) at (1.65,-0.20) {};
  \node[bspd] (sbo) at (3.75,0.20) {};
  \node[rspd] (sbi) at (3.75,-0.20) {};
  \draw[bluewire] (sao) -- (sf.160);
  \draw[redwire]  (sai) -- (sf.200);
  \draw[bluewire] (sf.20)  -- (sbo);
  \draw[redwire]  (sf.-20) -- (sbi);
  \node at (4.5,0) {$\leftarrow$};
  \node[cospanblue, minimum height=0.9cm] (sR) at (5.4,0) {};
  \node[bspd] at (5.4,0.20) {};
  \node[rspd] at (5.4,-0.20) {};
  \node[plabel, anchor=east] at ([shift={(-3pt,0.20cm)}]sL.west) {$(b,\mathrm{o},s)$};
  \node[plabel, anchor=east] at ([shift={(-3pt,-0.20cm)}]sL.west) {$(\bot,\mathrm{it},s)$};
  \node[plabel, anchor=west] at ([shift={(3pt,0.20cm)}]sR.east) {$(s,\mathrm{o},\top)$};
  \node[plabel, anchor=west] at ([shift={(3pt,-0.20cm)}]sR.east) {$(s,\mathrm{it},\top)$};
  \node at (7.4,0) {$\mapsto$};
  \node[morphismblue, minimum height=0.95cm] (sd) at (9.0,0) {$s$};
  \draw[bluewire] ($(sd.south west)!0.70!(sd.north west)$) -- ++(-0.55,0);
  \draw[redwire]  ($(sd.south west)!0.30!(sd.north west)$) -- ++(-0.55,0);
  \draw[bluewire] ($(sd.south east)!0.70!(sd.north east)$) -- ++(0.55,0);
  \draw[redwire]  ($(sd.south east)!0.30!(sd.north east)$) -- ++(0.55,0);

  \begin{scope}[yshift=-1.5cm]
    \node[font=\small] at (-2.15,0) {$g_{b,\mathrm{chk}}$};
    \node[cospanblue, minimum height=0.6cm] (bL) at (0,0) {};
    \node[bspd] at (0,0) {};
    \node at (0.9,0) {$\rightarrow$};
    \node[cospangrey, minimum height=1.0cm, minimum width=2.6cm] (bA) at (2.7,0) {};
    \node[mbox] (bf) at (2.7,0) {$b$};
    \node[bspd] (bao) at (1.65,0) {};
    \node[bspd] (bbo) at (3.75,0) {};
    \draw[bluewire] (bao) -- (bf.west);
    \draw[bluewire] (bf.east) -- (bbo);
    \node at (4.5,0) {$\leftarrow$};
    \node[cospanblue, minimum height=0.6cm] (bR) at (5.4,0) {};
    \node[bspd] at (5.4,0) {};
    \node[plabel, anchor=east] at ([shift={(-3pt,0cm)}]bL.west) {$(a,\mathrm{o},b)$};
    \node[plabel, anchor=west] at ([shift={(3pt,0cm)}]bR.east) {$(b,\mathrm{o},s)$};
    \node at (7.4,0) {$\mapsto$};
    \node[morphismblue, minimum size=0.75cm] (bsd) at (9.0,0) {$b$};
    \draw[bluewire] (bsd.west) -- ++(-0.55,0);
    \draw[bluewire] (bsd.east) -- ++(0.55,0);
  \end{scope}

  \begin{scope}[yshift=-3.0cm]
    \node[font=\small] at (-2.15,0) {$g_{b,\mathrm{skip}}$};
    \node[cospanblue, minimum height=0.6cm] (kL) at (0,0) {};
    \node[bspd] at (0,0) {};
    \node at (0.9,0) {$\rightarrow$};
    \node[cospangrey, minimum height=1.0cm, minimum width=2.6cm] (kA) at (2.7,0) {};
    \node[mbox] (kf) at (2.7,0) {$b$};
    \node[bspd] (kao) at (1.65,0) {};
    \node[bspd] (kbo) at (3.75,0) {};
    \draw[bluewire] (kao) -- (kf.west);
    \draw[bluewire] (kf.east) -- (kbo);
    \node at (4.5,0) {$\leftarrow$};
    \node[cospanblue, minimum height=0.6cm] (kR) at (5.4,0) {};
    \node[bspd] at (5.4,0) {};
    \node[plabel, anchor=east] at ([shift={(-3pt,0cm)}]kL.west) {$(\bot,\mathrm{o},b)$};
    \node[plabel, anchor=west] at ([shift={(3pt,0cm)}]kR.east) {$(b,\mathrm{o},s)$};
    \node at (7.4,0) {$\mapsto$};
    \node[morphismblue, minimum size=0.75cm] (ksd) at (9.0,0) {$b$};
    \draw[bluewire] (ksd.west) -- ++(-0.55,0);
    \draw[bluewire] (ksd.east) -- ++(0.55,0);
  \end{scope}
\end{tikzpicture}
\unskip
  \caption{The generators $g_s$, $g_{b,\mathrm{chk}}$ and $g_{b,\mathrm{skip}}$ of Fig.~\ref{fig:re-ocpn}, each as a decorated cospan
    (left, boundary $\to$ apex $\leftarrow$ boundary) and as the string diagram it denotes (right).
    $g_s$ is the object-centric sync. Two object types meet on the left boundary
    and leave on the right, joined by the apex hyperedge $s$. The silent $\tau$
    contributes no generator and gives $b$ two contexts, $g_{b,\mathrm{chk}}$
    (predecessor $a$) and $g_{b,\mathrm{skip}}$ (the open end $\bot$, traced
    through $\tau$), differing only in their left boundary port.}
  \label{fig:re-ocpn-cospans}
\end{figure*}

\begin{exmp}[A small OCPN and its generator cospans]
  \label{ex:running-pn-cospan}
  Take the object-centric Petri net of Fig.~\ref{fig:re-ocpn}, over the object
  types $\mathrm{order}$ and $\mathrm{item}$. Its observable transitions are
  $a,b,s$, with one silent transition $\tau$. The boundary nodes $\bot$ and $\top$ are not
  transitions of the net. They mark where objects enter and leave, and the walk ends there at the
  open ends of Definition~\ref{def:contexts-general}. Instantiating
  Theorem~\ref{thm:petri-to-cospan-signature} reads a generator cospan off
  each transition--context pair (Fig.~\ref{fig:re-ocpn-cospans}). Every edge has weight $1$,
  so each constraint system $\Lambda_{a,c}$ fixes $n_p=1$ on every leg and each wire carries a
  single object. Write $(x,\mathrm{o},y)$ for $(x,\{\mathrm{o}\},y)$.

  \emph{The check $a$.} The check has one observable predecessor route and one
  observable successor, so it has a single context and a single generator,
  \[
    g_a = \Bigl(\{(\bot,\mathrm{o},a)\}
      \xrightarrow{\iota} A \xleftarrow{\iota}
      \{(a,\mathrm{o},b)\},\; d_a\Bigr),
  \]
  with one order wire on each boundary.

  \emph{The sync $s$.} The transition $s$ consumes the registered order together
  with the item and emits both onward, so its generator
  \begin{gather*}
    g_s = \Bigl(\{(b,\mathrm{o},s),(\bot,\mathrm{it},s)\}
      \xrightarrow{\iota} A \xleftarrow{\iota}\twocolbreak 
      \{(s,\mathrm{o},\top),(s,\mathrm{it},\top)\},\; d_s\Bigr)
  \end{gather*}
  carries two typed wires on each boundary and the single apex hyperedge labelled
  $s$.

  \emph{The silent skip.} The check $a$ may be skipped through $\tau$. A silent
  transition is a transparent mediator (Remark~\ref{rem:ocpn-silent-typing}) and
  produces no generator. It gives $b$ a second context instead. Tracing the order
  thread back from $b$ reaches either $a$ (the checked route) or, through $\tau$,
  directly the open end $\bot$ (the skip route), so $b$ has two generator cospans,
  \begin{gather*}
    g_{b,\mathrm{chk}} : \{(a,\mathrm{o},b)\}\to\{(b,\mathrm{o},s)\},\qquad\twocolbreak
    g_{b,\mathrm{skip}} : \{(\bot,\mathrm{o},b)\}\to\{(b,\mathrm{o},s)\},
  \end{gather*}
  differing only in their left boundary port. This is the Petri-net form of an
  exclusive choice. The two contexts give one generator per route, re-identified at the shared
  label $b$ by the construction above.
\end{exmp}

The example shows that the split place and its silent $\tau$
introduce no structural difference in the signature, as
Remark~\ref{rem:ab-place-vs-tau} records and the internal-routing quotient of
Corollary~\ref{cor:internal-routing-quotient} makes precise in the conversion
certificate.

\begin{figure*}[tp]
  \centering
  \begin{tikzpicture}[inline, scale=1.0,
    cut/.style={dashed, thick, gray!70},
    cutlab/.style={font=\footnotesize, gray!80}]

  \node[morphismblue, minimum height=1.8cm] (g1) at (0,0)   {$\gamma_1$};
  \node[morphismblue] (a)  at (2.2,0.45) {$a$};
  \node[morphismblue] (b)  at (3.9,0.45) {$b$};
  \node[morphismblue, minimum height=1.8cm] (s)  at (6.0,0) {$s$};
  \node[morphismblue, minimum height=1.8cm] (g2) at (7.8,0) {$\gamma_2$};

  \draw[bluewire] ($(g1.north east)!0.25!(g1.south east)$) -- (a.west);
  \draw[bluewire] (a.east) -- (b.west);
  \draw[bluewire] (b.east) -- ($(s.north west)!0.25!(s.south west)$);
  \draw[bluewire] ($(s.north east)!0.25!(s.south east)$) -- ($(g2.north west)!0.25!(g2.south west)$);
  \draw[redwire]  ($(g1.north east)!0.75!(g1.south east)$) -- ($(s.north west)!0.75!(s.south west)$);
  \draw[redwire]  ($(s.north east)!0.75!(s.south east)$) -- ($(g2.north west)!0.75!(g2.south west)$);

  \draw[cut] (3.05,-0.7) -- (3.05,0.7);
  \node[cutlab, above] at (3.05,0.7) {$M_1$};
  \draw[cut] (4.95,-0.7) -- (4.95,0.7);
  \node[cutlab, above] at (4.95,0.7) {$M_2$};

  \node[cutlab, align=center] at (3.9,-1.5)
    {each cut is a marking; one crossed wire $=$ one token of that type\\
     $M_1=M_2=\{1\ \textcolor{bluecol}{\mathrm{order}},\ 1\ \textcolor{redcol}{\mathrm{item}}\}$};
\end{tikzpicture}
\unskip
  \caption{Markings as cuts. A vertical cut through the connected string diagram of
    the net of Fig.~\ref{fig:re-ocpn} fixes a marking. Each wire it crosses is one
    token of that wire's type. The cuts $M_1$ (after $a$) and $M_2$ (after $b$) each
    hold one order and one item token, the contents of the intervening places.
    The diagram is closed by a start generator $\gamma_1$ and an end generator $\gamma_2$, a
    choice of the analysis (Section~\ref{subsec:analysis}), and sweeping the cut from $\gamma_1$
    to $\gamma_2$ replays the net.}
  \label{fig:re-pn-slice}
\end{figure*}

\begin{remark}[Markings as cuts]\label{rem:pn-marking-cut}
  A marking of the net corresponds to a vertical cut through its connected string
  diagram. The wires the cut crosses are exactly the tokens present, one per wire and
  typed by the wire (Fig.~\ref{fig:re-pn-slice}). Sweeping the cut from $\gamma_1$ to
  $\gamma_2$ reads off one execution as a sequence of markings, and the number of
  crossings of a given colour is the token count of that object type, recovering the
  marking semantics of Section~\ref{sec:hypergraphs} on this example. This is the
  Petri-net counterpart of the causal-net slicing of Section~\ref{sec:causal-nets}.
\end{remark}

